% Width: 169 mm (full width, figure*); height: 103 mm.
% Use at native size with \input; all labels are at least footnotesize.
% Population schematic, not random samples or an ICA optimization result.
% Independent sources s_1,s_2 ~ Uniform[-sqrt(3),sqrt(3)] have mean 0, Var=1.
% R(theta)=[cos(theta),-sin(theta);sin(theta),cos(theta)].
% x=R(30 deg)s, Cov(x)=I_2; z=R(-30 deg)x=s.
% The square outline is the support of the continuous population.
% Dots are the deterministic Cartesian grid {(j/2,k/2):j,k=-3,...,3}.
% Its 49 equally weighted points also have mean 0 and covariance I_2
% with denominator n, but do not constitute a random sample.
% Plot scale: 11 mm per dimensionless coordinate unit in both panels.
% Matrix convention: S,X,Z in R^{49 x 2} store samples as rows.
% X=S R(30 deg)^T, W=R(30 deg), Z=X W=S.
% Dashed arrows are columns of W; PCA has no preferred direction because
% its population eigenvalues are equal, even though x_1,x_2 are dependent.
% ICA only identifies these non-Gaussian sources up to scale and permutation.
\begin{tikzpicture}[
  x=1mm,y=1mm,font=\footnotesize,text=BookInk,
  ica axis/.style={-{Stealth[length=1.7mm]},draw=BookMuted,line width=.6pt},
  ica direction/.style={-{Stealth[length=1.8mm]},draw=FigureModel,
    line width=.8pt,dashed}
]
  \path[use as bounding box] (0,0) rectangle (169,103);
  \node[font=\small] at (38,98) {（a）白化混合：不相关};
  \node[font=\small] at (131,98) {（b）源坐标：独立};
  \node at (38,89) {$\bm{x}=\bm{R}(30^\circ)\bm{s}$};
  \node at (131,89) {$\bm{z}=\bm{R}(-30^\circ)\bm{x}=\bm{s}$};

  \begin{scope}[shift={(38,52)},x=11mm,y=11mm]
    \begin{scope}[rotate=30]
      \draw[draw=FigureInput,fill=FigureInput!8!BookPaper,line width=.8pt]
        ({-sqrt(3)},{-sqrt(3)}) rectangle ({sqrt(3)},{sqrt(3)});
      \foreach \j in {-3,...,3}{
        \foreach \k in {-3,...,3}{
          \fill[FigureInput] ({\j/2},{\k/2}) circle (1.1pt);
        }
      }
    \end{scope}
    \draw[ica axis] (-2.65,0) -- (2.65,0) node[right] {$x_1$};
    \draw[ica axis] (0,-2.65) -- (0,2.65) node[above] {$x_2$};
    \foreach \tick in {-2,2}{
      \draw[BookMuted] (\tick,-.06) -- (\tick,.06);
      \node[below,inner sep=1.5pt] at (\tick,-.06) {$\tick$};
      \draw[BookMuted] (-.06,\tick) -- (.06,\tick);
      \node[left,inner sep=1.5pt] at (-.06,\tick) {$\tick$};
    }
    \node[below left,inner sep=1.5pt] at (0,0) {$0$};
    \begin{scope}[rotate=30]
      \draw[ica direction] (-2.2,0) -- (2.4,0)
        node[above right,inner sep=1.5pt] {$\bm{w}_1$};
      \draw[ica direction] (0,-2.2) -- (0,2.4)
        node[above left,inner sep=1.5pt] {$\bm{w}_2$};
    \end{scope}
  \end{scope}

  \draw[book arrow,draw=FigureModel] (76,52) -- (96,52);
  \node[align=center] at (86,63) {ICA解混\\旋转$-30^\circ$};
  \node[align=center] at (86,42) {$\bm{z}=\bm{W}^{\top}\bm{x}$\\$\bm{W}=\bm{R}(30^\circ)$};

  \begin{scope}[shift={(131,52)},x=11mm,y=11mm]
    \draw[draw=FigureOutput,fill=FigureOutput!8!BookPaper,line width=.8pt]
      ({-sqrt(3)},{-sqrt(3)}) rectangle ({sqrt(3)},{sqrt(3)});
    \foreach \j in {-3,...,3}{
      \foreach \k in {-3,...,3}{
        \fill[FigureOutput] ({\j/2},{\k/2}) circle (1.1pt);
      }
    }
    \draw[ica axis] (-2.65,0) -- (2.65,0)
      node[right,align=center,inner sep=1pt] {$z_1$\\$=s_1$};
    \draw[ica axis] (0,-2.65) -- (0,2.65) node[above] {$z_2=s_2$};
    \foreach \tick in {-2,2}{
      \draw[BookMuted] (\tick,-.06) -- (\tick,.06);
      \node[below,inner sep=1.5pt] at (\tick,-.06) {$\tick$};
      \draw[BookMuted] (-.06,\tick) -- (.06,\tick);
      \node[left,inner sep=1.5pt] at (-.06,\tick) {$\tick$};
    }
    \node[below left,inner sep=1.5pt] at (0,0) {$0$};
  \end{scope}

  \node[align=center] at (38,15)
    {$\operatorname{Cov}(\bm{x})=\bm{I}_2$\\PCA仅凭二阶量无方向偏好};
  \node[align=center] at (131,15)
    {$\operatorname{Cov}(\bm{z})=\bm{I}_2$\\此均匀分布在源坐标下分解为独立分量};
  \node at (84.5,3)
    {$s_1,s_2$独立且均服从$\mathcal{U}[-\sqrt{3},\sqrt{3}]$；坐标无量纲，格点只示意总体支持域};
\end{tikzpicture}%
